\section{Introduction}\label{sec:intro}

A sovereign that borrows in its own reserve currency, and whose central bank can buy its bonds, does not face a debt wall: default is a choice, not a constraint. When interest costs rise and the budget does not respond, the gap is still paid, by one of three groups. Taxpayers or program beneficiaries pay if fiscal policy adjusts. Holders of nominal claims pay if the central bank tolerates the inflation that erodes the government's nominal liabilities. And taxpayers pay again if the central bank raises real rates instead: the debt that reprices then costs more, and the inflation transfer is taken back.

This paper prices these routes. For the United States from 1980 to 2025, and for the United Kingdom and Japan, it measures how much fiscal adjustment a rate shock requires, how much inflation would substitute for that adjustment, and how much monetary tightening the substitution can survive. The answers to the last two questions are set by the consolidated balance sheet of the Treasury and the central bank, and above all by how the central bank funds the bonds it holds.

\paragraph{Two layers.} In a transparent model of consolidated debt dynamics, the market rate rises with debt, and the average rate on outstanding debt reprices toward it at a speed set by maturity. Debt is locally stable if and only if
\begin{equation}
(1-\varphi)\,(r+\psi b)<g,
\end{equation}
where $\varphi$ is the share of marginal interest cost offset by the primary surplus. The \emph{fiscal layer} is maturity-free: the repricing speed cancels from this condition for any maturity structure, so the minimum offset $\varphi^*=1-g/(r+\psi b)$ depends on $r-g$, on the debt sensitivity of rates $\psi$ and on debt $b$, not on maturity.

The \emph{inflation layer} is where maturity and the central bank matter. Sustained surprise inflation erodes only nominal liabilities that have not yet repriced, and currency, which never reprices. The inflation that covers a missing offset $u$ after a permanent rate rise $\Delta r$ is
\begin{equation}
\Delta\pi(H)=u\,\Delta r\,R,\qquad R=\frac{\int_0^H P(h)\,dh}{\int_0^H E(h)\,dh},
\end{equation}
where $P(h)$, the \emph{rollover clock}, is the share of the rate shock that has reached the average rate after $h$ years, and $E(h)$ is the share of the consolidated balance sheet that inflation can still erode. $R$ is the price of the inflation route. If the central bank raises the real rate by $\kappa$ per point of inflation, repriced debt pays more and the requirement becomes $u\,\Delta r\,R/(1-\kappa R)$. It is infinite beyond a \emph{monetary tolerance threshold} $\kappa^*=1/R$. The same ratio therefore sets both the price of the inflation route and how much monetary tightening the route can survive.

\paragraph{Measurement and tests.} We build the U.S. clock security by security from the Monthly Statement of the Public Debt for 2001--2025. We net the Fed's holdings CUSIP by CUSIP, add reserves and reverse repos as overnight liabilities, and add currency to the erosion base. Treasury Bulletin tables extend the series to 1980, and the reconstruction matches official totals exactly. Three tests check it. Over 22 year-end origins, it predicts the Treasury's average interest rate a year ahead within 0.05pp, a sixth of the error of a clock based on weighted-average maturity. Frozen at end-2021, it tracks the 1.9pp rise in the average rate through 2025 and, with reserves repricing overnight, reproduces the Fed's cumulative operating loss: a deferred asset of $-\$245$bn predicted, against $-\$243$bn actual. And from the end-2020 balance sheet, the 2021--23 inflation surprise transferred about 11\% of GDP from holders of nominal liabilities by 2023. That is 1.8 times what full Fisher repricing implies, because new debt was priced for far less inflation than followed; higher real rates then took most of it back.

\paragraph{Findings for the United States.} The fiscal threshold is ordinary. At about 0.47 in 2023--25, $\varphi^*$ is within its 1980--2007 range; the exception was 2008--21, when it was negative. The price of falling short is not. At the end-2025 clock, each 0.1 of missing offset after a permanent 1pp rate rise costs about 0.2pp of extra inflation a year for a decade. With no offset the cost is 1.0--1.1pp a year, the highest since at least 1980. This ranking holds for any common fiscal response below about 0.35 and in 99\% of parameter draws, provided today's $\psi$ is at least about 2bp. Holding debt/GDP stable at CBO's baseline deficits would take about 4pp a year.

QE raised the price. Before 2008 the Fed held Treasuries against currency, which inflation erodes, and consolidation roughly halved the price; since 2009 reserves that reprice overnight fund its book, and consolidation no longer lowers it. QE's reserve funding accounts for 50--80\% of the rise in the consolidated ratio since 2007. QE also lowered the monetary tolerance threshold, by about two fifths at five-, ten- and fifteen-year horizons. Over ten years, $\kappa^*$ fell from 0.64--0.82 before 2008 to 0.39--0.50 in every year since 2009, below the 0.5 of a \citet{taylor1993} rule. Beyond a few years, the inflation route is open only to a central bank that responds to inflation much less than its usual rule.

\paragraph{Three sovereigns.} The same measures for the United Kingdom (2007--2025) and Japan (fiscal years 2020--2024) show that these are properties of the consolidated balance sheet, not of debt management. The UK issues the longest debt of the three, and its own-debt price is about 0.9, against 2.1--2.7 in the United States; yet at the 2021 QE peak its consolidated price, 2.47, matched the U.S. 2.46. Japan's zero-rate reserve tier kept its price near 0.8 until the March 2024 reform raised it to about 2. The post-2020 inflations show two of the three routes. In the United States and the UK, higher real rates took most of the transfer back, in the UK almost entirely through interest on reserves (5.7\% of GDP). In Japan, where yields were held down, holders of nominal claims bore about 18\% of GDP.

\paragraph{The fiscal response.} Whether the threshold is met is a policy choice, and the paper does not need to estimate it: every result is reported as a function of the response. The best available evidence suggests the response is weak. \citet{auerbach2024} estimate a legislated U.S. offset of higher interest costs of about 0.39 before 2004 and about zero since. Our break tests confirm a change but date it to the mid-1990s and 2009--13, and UK fiscal events, where market-driven revisions to debt interest identify the response, give about 0.15 (Appendix~\ref{app:fiscal}).

\paragraph{Contribution.} None of the ingredients is new on its own. Debt dynamics with endogenous rates come from \citet{mian2025} and \citet{lorenzoni2019}; consolidated Treasury--Fed maturity from \citet{greenwood2014}, the \citet{tbac2020,tbac2026} and the \citet{obr2021}; maturity and the inflation tax from \citet{cochrane2001,cochrane2022}, \citet{hilscher2022} and \citet{barro2026}; and fiscal responses from \citet{bohn1998} and \citet{auerbach2024}. The paper adds three things. It proves that the fiscal threshold is maturity-free, and derives closed forms for the price of the inflation route and for the monetary tolerance threshold in terms of the measured clock. It builds and tests a 1980--2025 series of the full consolidated repricing profile of the United States, which we have not found elsewhere, and the same measures for the UK and Japan. And it quantifies the effect of consolidation on the inflation requirement, which \citet{barro2026} note would be desirable but for which ``data are not available.''

\paragraph{What the results rest on.} The results differ in how much they depend on assumptions. The measurement and the accounting are the most secure; the levels of the requirement are orders of magnitude; the level of the fiscal response is the least secure, and the paper does not rely on it.

\begin{center}
\footnotesize
\begin{tabularx}{\textwidth}{>{\raggedright\arraybackslash}X>{\raggedright\arraybackslash}X>{\raggedright\arraybackslash}X>{\raggedright\arraybackslash}X}
\toprule
Result & Basis & Depends on & Reliability \\
\midrule
The rollover clock and the consolidated stock & Security-level data, exact against official totals; backtest error 0.05 pp; the Fed's 2022–25 loss reproduced & Bucket correction for 1980–2002 & High \\
The fiscal threshold does not depend on maturity & Proof (Section 3.2) & Linear debt dynamics; constant $\psi$ & High, within the model class \\
Funding central-bank bonds with reserves raises the price of the inflation route & Accounting; the same in the United States, the UK and Japan; Japan's 2024 reform & Ignores the return on the Fed's MBS & High \\
QE lowered the monetary tolerance threshold $\kappa^*$ & Same accounting, three countries & First order; the response lasts as long as the surprise & High for the fall (about two fifths at H = 5, 10, 15); below 0.5 at H = 10 and 15, not at H = 5 \\
Markets priced far less of the post-2020 surprises than full Fisher repricing & U.S. and UK tests (1.8 times in both); Japan & UK and Japan: closed portfolio & High \\
2023–25 has the largest requirement since 1980 & Year-by-year accounting & Today's $\psi$ at least about 2bp (3bp at $\hat\varphi$ = 0.25) & Medium \\
Levels of the requirement (1.0–1.1 pp; about 4 pp to hold debt/GDP) & Same & $\psi$, g, H; currency demand held fixed & Order of magnitude \\
The level of the fiscal response $\hat\varphi$ & UK fiscal events about 0.15; U.S. not identified & — & Low; not needed for the results \\
\bottomrule
\end{tabularx}
\end{center}


\paragraph{Plan of the paper.} Section~\ref{sec:literature} reviews related work and Section~\ref{sec:model} presents the model. Section~\ref{sec:measurement} builds and tests the clock, and Section~\ref{sec:pricing} prices the fiscal gap from 1980 to 2025. Section~\ref{sec:international} turns to the United Kingdom and Japan, Section~\ref{sec:policy} to policy options, and Section~\ref{sec:conclusion} concludes. Appendix~\ref{app:timing} treats timing (a one-time price jump and a ceiling on the surplus), and Appendix~\ref{app:fiscal} the evidence on the fiscal response.
